\subsection{THE AXIOMS}

A logical theory is any theory \(\mathfrak{G}\) in which the schemes S1 to S4 below provide implicit axioms.

S1. If \(A\) is a relation in \(\mathfrak{C}\) , the relation \((A \text{ or } A) \Longrightarrow A\) is an axiom of \(\mathfrak{C}\)\footnote{This scheme may be expressed without using the letter A or the abbreviating symbol \(\Longrightarrow\) as follows: whenever we have a relation, we obtain a theorem by writing, from left to right, \(\vee\), \(\neg\), \(\vee\), and then the given relation three times. The reader may, as an exercise, translate in a similar way the expressions of the other schemes.}.

S2. If A and B are relations in C, the relation \(A \Longrightarrow (A \text{ or } B)\) is an axiom of C.

S3. If A and B are relations in \(\mathfrak{G}\) , the relation (A or B) \(\Longrightarrow\) (B or A) is an axiom of G.

S4. If A, B, and C are relations in \(\mathfrak{C}\) , the relation

\[
(A \Rightarrow B) \Rightarrow ((C \text {   or   } A) \Rightarrow (C \text {   or   } B))
\]

is an axiom of \(\mathfrak{C}\) .

These rules are in fact schemes; let us verify this, for example for S2. Let R be a relation obtained by applying S2; then there are relations A and B in C such that R is the relation \(A \Longrightarrow (A \text{ or } B)\) . Let T be a term in C, let x be a letter, and let \(A'\) and \(B'\) be the relations \((T|x)A\) and \((T|x)B\) ; then \((T|x)R\) is the same as \(A' \Longrightarrow (A' \text{ or } B')\) , and can therefore be obtained by applying S2.

Intuitively, the rules S1 through S4 merely express the meaning which is attached to the words ``or'' and ``implies'' in the usual language of mathematics\footnote{In everyday speech, the word ``or'' has two different meanings, according to the context: when we link two statements by the word ``or'' we may mean to assert at least one of the two (and possibly both together), or we may mean to assert one to the exclusion of the other.}.

If a logical theory \(\mathfrak{C}\) is contradictory, every relation in \(\mathfrak{C}\) is a theorem in \(\mathfrak{C}\) . For let \(A\) be a relation in \(\mathfrak{C}\) such that \(A\) and ``not \(A\)'' are theorems in \(\mathfrak{C}\) , and let \(B\) be any relation in \(\mathfrak{C}\) . By S2, (not \(A\) ) \(\Longrightarrow\) ((not \(A\) ) or \(B\) ) is a theorem in \(\mathfrak{C}\) ; therefore, by C1 ( \(\S 2\) , no. 2), ``(not \(A\) ) or \(B\)'', that is to say \(A\Longrightarrow B\) , is a theorem in \(\mathfrak{C}\) . A second application of C1 shows that \(B\) is a theorem in \(\mathfrak{C}\) .

From now on \(\mathfrak{C}\) will denote a logical theory.

